\family{Moraceae}{Fig family}
\label{fam:moraceae}

Trees, shrubs or climbers with white latex. Leaves alternate, simple,
with stipules that fall early and leave a scar round the twig. Flowers
small and unisexual, crowded in heads, spikes or catkins, or enclosed
in a fleshy receptacle. Fruits often fused into a fleshy compound
fruit. Two species in the Arami Uplands, both timber trees.

\section*{Key to the species}
\begin{florakey}
	\lead{1a}{Leaves unequal at the base, rough; male flowers in small,
		flat heads; fruit a single red, velvety drupe}{\taxon{Antiaris
			toxicaria}}
	\lead{1b}{Leaves equal at the base, smooth; flowers in catkins, the
		male catkins long and pendent; fruit a fleshy green compound fruit
		like a large mulberry}{\taxon{Milicia excelsa}}
\end{florakey}

\species{Antiaris toxicaria}{Lesch.}
\label{sp:MOR01}

\field{Synonyms} \sciname{Antiaris africana}{Engl.}

\field{Local and trade names} Trade: Ako.

\field{Status} Native.

\field{Description} Deciduous tree to 40~m with small buttresses. Bark
smooth, grey to whitish, exuding abundant white latex when cut. Leaves
oblong-elliptic, 8--20~cm long, unequal at the base and rough to the
touch, especially on young trees. Male flowers in small, flat heads;
female flowers solitary. Fruit a red, velvety drupe about 2~cm long.

\field{Phenology} Flowering November to February; fruiting February
to April.

\field{Habitat} Lowland semi-deciduous forest.

\field{Notes} West African plants are referred to
\taxon{Antiaris toxicaria} subsp. \taxon{welwitschii}
\authority{(Engl.) C.C.Berg}.

\field{Uses} A light timber used for plywood and boxes. The inner bark
has been beaten into bark cloth \citep{hawthorne2006,louppe2008}.

\field{Specimens examined} Okeri (C4), 420~m, fl., 26~Feb 1935, \specimen{Harding}{1115}{OKF};
Igbo-Oba Forest Reserve (C3), 470~m, fr., 7~Mar 1936, \specimen{Harding}{1233}{OKF};
Ireke Ferry (D4), 360~m, fr., 3~Mar 1938, \specimen{Harding}{1596}{OKF};
Ireke Ferry (D4), 360~m, st., 3~May 2019, \specimen{Alade}{112}{AUH};
Olowo Stream (D3), 440~m, fl., 11~Dec 2019, \specimen{Alade}{134}{AUH}.

\field{Distribution} Grid squares C3, C4, D3, D4.

\begin{center}
	\distmap{C/3, C/4, D/3, D/4}
\end{center}

\species{Milicia excelsa}{(Welw.) C.C.Berg}
\label{sp:MOR02}

\field{Synonyms} \sciname{Morus excelsa}{Welw.}

\field{Local and trade names} Yoruba: Iroko. Trade: Iroko.

\field{Status} Native.

\field{Description} Deciduous tree to 50~m with a long, clear bole and
small buttresses. Bark pale grey, scaly, exuding white latex. Leaves
elliptic, 10--20~cm long, pointed at the tip and rounded at the base.
Male and female flowers on separate trees; male catkins 10--20~cm
long, pendent and white; female catkins shorter and green. Fruit a
fleshy, green compound fruit 5--7~cm long, like a large mulberry.

\field{Phenology} Flowering October to November; fruiting December to
February.

\field{Habitat} Lowland semi-deciduous forest and forest--savanna
mosaic; large trees are often left standing in farmland.

\field{Notes} This species was long placed in \taxon{Chlorophora};
\citet{berg1982} reinstated the genus \taxon{Milicia} for it.

\field{Uses} One of the most important timbers of West Africa, very
durable and used for building, furniture, flooring and boats. Old
trees are held sacred in many communities \citep{keay1989,louppe2008}.

\field{Specimens examined} Okeri (C4), 420~m, fr., 11~Dec 1936, \specimen{Harding}{1313}{OKF};
Agbede (C2), 540~m, fr., 11~Dec 1937, \specimen{Harding}{1511}{OKF};
Aramiwo Foot (B2), 640~m, fr., 4~Feb 1962, \specimen{Olaniyan}{325}{SVN};
Ewe Stream gallery forest (B3), 580~m, fr., 21~Jan 1967, \specimen{Olaniyan}{823}{SVN};
Oke-Asa farmland (B4), 690~m, fr., 11~Dec 2022, \specimen{Alade}{527}{AUH}.

\field{Distribution} Grid squares B2, B3, B4, C2, C4.

\begin{center}
	\distmap{B/2, B/3, B/4, C/2, C/4}
\end{center}